\documentclass[11pt]{article}
\usepackage{amsmath,amssymb,amsthm,bm}
\usepackage[margin=1in]{geometry}

\newcommand{\Sig}{\Sigma_e}
\newcommand{\E}{\mathbb{E}}
\newcommand{\Var}{\mathrm{Var}}
\newcommand{\diag}{\mathrm{diag}}
\newcommand{\Rk}{R_k^{\mathrm{bag},\infty}}
\newcommand{\Qsub}{Q_k^{\mathrm{sub}}}
\newcommand{\Qopt}{Q^{\mathrm{opt}}}
\newcommand{\iotab}{\bm{\iota}}

\newtheorem{assumption}{Assumption}
\newtheorem{theorem}{Theorem}
\newtheorem{lemma}{Lemma}
\newtheorem{proposition}{Proposition}
\newtheorem{corollary}{Corollary}
\newtheorem{remark}{Remark}

\title{The Infinite-Bagged RSM Remainder $R_k$:\\
What Holds Exactly, and What Needs Assumptions}
\date{June 19, 2026}

\begin{document}
\maketitle

\begin{abstract}
This note separates, as cleanly as possible, the results on the infinite-bagged
RSM excess risk $\Rk=Q_k^{\mathrm{bag},\infty}-\Qopt$ into two tiers: (i) what is
\emph{exact} --- true under the common-loading structure with no further assumption,
no approximation, and no asymptotics; and (ii) what \emph{rate} statements
($\Rk\asymp z_k^2$, hence the cube-root rule $k^\ast\asymp T^{1/3}$) require, namely
two explicit, checkable, and minimal assumptions. The exact tier gives a computable
two-sided bracket and a complete characterization of the optimal $k$; the rate tier
is governed entirely by which of two assumptions is binding, mapping the effective
exponent $\alpha\in[1,2]$ to the bounded-coherence/heterogeneity regime.
\end{abstract}

%==========================================================================
\section{Setup and notation}
Assume the common-loading covariance structure of the forecast errors,
\[
  \Sig = q\,\iotab\iotab' + D,\qquad D=\diag(\sigma_1^2,\dots,\sigma_m^2),
  \qquad \tau_i=\sigma_i^{-2}>0,
\]
where $q\ge0$ is the variance of the common (unforecastable) shock and $\tau_i$ is the
precision of forecaster $i$'s idiosyncratic error. Write
\[
  A_m=\sum_{i=1}^m\tau_i=m\bar\tau,\qquad \ell_i=\frac{\tau_i}{A_m},
  \qquad \sum_{i=1}^m\ell_i=1,
\]
and for a subset $S$ of size $k$, $A_S=\sum_{j\in S}\tau_j$. RSM draws subsets $S$
uniformly at random, fits the constrained (sum-to-one) optimal weights
$v_{S,i}=\mathbf 1\{i\in S\}\,\tau_i/A_S$ within each, and bags. The
infinite-bagging ($G\to\infty$) weight is $\bar v_{k,i}=\E_S[v_{S,i}]$, and the object
of interest is the population excess risk
\[
  \Rk=Q_k^{\mathrm{bag},\infty}-\Qopt=\bar v_k'\Sig\bar v_k-\Qopt .
\]
Two population functionals carry the entire story:
\[
  V_\delta=\frac1m\sum_{i=1}^m\Big(\frac{\tau_i}{\bar\tau}-1\Big)^2
  =\frac{\Var(\tau)}{\bar\tau^2}\quad(\text{normalized precision dispersion}),
  \qquad
  H=\overline{\sigma^2}\,\bar\tau=\Big(\tfrac1m\textstyle\sum_i\sigma_i^2\Big)\bar\tau\ge1,
\]
the latter being the arithmetic-vs-harmonic gap (Cauchy--Schwarz), with $H=1$ iff all
precisions are equal. Throughout, $z_k=\dfrac1k-\dfrac1m=\dfrac{m-k}{km}$.

%==========================================================================
\section{Tier I: what holds exactly (no assumptions)}
Every statement in this section is an identity or inequality that holds for an
arbitrary positive precision vector $(\tau_1,\dots,\tau_m)$ and every $1\le k\le m$.
No balanced-leverage, no distributional assumption, no Taylor expansion, no
$T\to\infty$ or $m\to\infty$ limit.

\begin{proposition}[Oracle weights are leverage scores]\label{prop:oracle}
The constrained-optimal full combination is $w_i^\ast=\tau_i/A_m=\ell_i$, with
$\Qopt=q+1/A_m$, and for any subset $Q(S)=q+1/A_S$.
\end{proposition}
\begin{proof}
Since $\iotab'w=1$, $w'q\iotab\iotab'w=q$ is constant, so minimizing $w'\Sig w$ reduces
to minimizing $\sum_i w_i^2\sigma_i^2$ subject to $\sum_iw_i=1$; the Lagrange solution is
$w_i\propto\tau_i$. Substituting gives $\Qopt=q+\sum_i\ell_i^2/\tau_i=q+1/A_m$. The subset
claim is the same argument restricted to $S$.
\end{proof}

\begin{proposition}[Exact quadratic-form identity]\label{prop:identity}
Because $w^\ast$ is the constrained minimizer, the cross term vanishes and
\[
  \Rk=(\bar v_k-w^\ast)'\Sig(\bar v_k-w^\ast)
  =\frac{1}{A_m}\sum_{i=1}^m\ell_i\Big(\frac{\bar v_{k,i}}{\ell_i}-1\Big)^2
  =\frac{1}{m^2}\sum_{i=1}^m\tau_i\Big(g_i-\frac1{\bar\tau}\Big)^2,
\]
where $g_i(k)=k\,\E_S[1/A_S\mid i\in S]$ and $\bar v_{k,i}=\tau_i g_i(k)/m$.
\end{proposition}
\begin{proof}
$\Sig w^\ast=(q+1/A_m)\iotab=\Qopt\iotab$, and since $\iotab'(\bar v_k-w^\ast)=0$
(both weight vectors sum to one), $(\bar v_k-w^\ast)'\Sig w^\ast=0$. The common
factor $q\iotab\iotab'$ then drops for the same reason, leaving
$\sum_i(\bar v_{k,i}-\ell_i)^2/\tau_i$; the two displayed forms are algebraic rewrites
using $w_i^\ast=\ell_i$ and $\tau_i=A_m\ell_i$.
\end{proof}

\begin{proposition}[Endpoints: RSM bridges averaging and the oracle]\label{prop:endpoints}
At $k=1$, $\bar v_1=\tfrac1m\iotab$ (simple averaging), so $R_1$ equals exactly the
excess risk of the equal-weight combination. At $k=m$, $\bar v_m=w^\ast$, so $R_m=0$.
Thus $k$ traces, in population, a path from simple averaging to the oracle combination.
\end{proposition}

\begin{lemma}[Two exact moment identities]\label{lem:moments}
For all $k$,
\[
  \sum_{i=1}^m\tau_i g_i=m\quad(\text{conservation law}),\qquad
  \sum_{i=1}^m g_i=km\,\E_S[1/A_S].
\]
\end{lemma}

\begin{proposition}[Jensen sandwich]\label{prop:jensen}
For all $1\le k\le m$,
\[
  0\le \Rk\le \Qsub-\Qopt=\E_S[1/A_S]-\frac1{A_m}.
\]
\end{proposition}
\begin{proof}
Lower bound: $w^\ast$ is the global sum-to-one minimizer and $\bar v_k$ is feasible.
Upper bound: $Q_k^{\mathrm{bag},\infty}=(\E_S v_S)'\Sig(\E_S v_S)\le\E_S[v_S'\Sig v_S]=\Qsub$
by Jensen ($\Sig\succeq0$). Subtract $\Qopt$ and use $Q(S)=q+1/A_S$, $\Qopt=q+1/A_m$.
\end{proof}

\begin{proposition}[Exact two-sided bracket]\label{prop:bracket}
For all $1\le k\le m$,
\[
  \frac{\big(k\,\E_S[1/A_S]-1/\bar\tau\big)^2}{\sum_i\sigma_i^2-m/\bar\tau}
  \;\le\;\Rk\;\le\;\E_S[1/A_S]-\frac1{A_m}.
\]
Both ends are computable exactly from $\{\tau_i\}$ and the scalar $\E_S[1/A_S]$. The
denominator $\sum_i\sigma_i^2-m/\bar\tau\ge0$ by Cauchy--Schwarz, with equality iff all
$\tau_i$ are equal (in which case $\Rk=0$). The lower bound is \emph{exact} whenever the
precision profile takes at most two distinct values, and a convergent hierarchy of exact
lower bounds follows by appending higher moment identities to Lemma~\ref{lem:moments}.
\end{proposition}
\begin{proof}
The upper bound is Proposition~\ref{prop:jensen}. For the lower bound, minimize
$\sum_i\tau_ig_i^2$ over all vectors satisfying the two identities of
Lemma~\ref{lem:moments}. In the inner product $\langle x,y\rangle_\tau=\sum_i\tau_ix_iy_i$
these read $\langle g,\iotab\rangle_\tau=m$ and $\langle g,c\rangle_\tau=km\,\E_S[1/A_S]$
with $c_i=1/\tau_i$. With $\|\iotab\|_\tau^2=A_m$, $\langle c,\iotab\rangle_\tau=m$, and
$\|c\|_\tau^2=\sum_i1/\tau_i$,
\[
  \min\sum_i\tau_ig_i^2=\frac{m}{\bar\tau}
  +\frac{m^2\big(k\E_S[1/A_S]-1/\bar\tau\big)^2}{\sum_i1/\tau_i-m/\bar\tau}.
\]
Substituting into Proposition~\ref{prop:identity} and using
$\tfrac1{m^2}\cdot\tfrac{m}{\bar\tau}=1/A_m$, the oracle term cancels.
\end{proof}

\begin{corollary}[The unconditional open bracket]\label{cor:open}
Expanding both ends of Proposition~\ref{prop:bracket} to leading order, using the exact
moments $\E_S[A_S]=k\bar\tau$, $\Var_S(A_S)=k\frac{m-k}{m(m-1)}V$ with
$V=\sum_i(\tau_i-\bar\tau)^2=m\bar\tau^2V_\delta$, and
$\E_S[1/A_S]=\frac1{k\bar\tau}+\frac{\Var_S(A_S)}{(k\bar\tau)^3}+\cdots$,
\[
  \underbrace{\frac{V_\delta^2}{m(H-1)\,\bar\tau}\,z_k^2}_{\text{floor}}
  \;\lesssim\;\Rk\;\lesssim\;
  \underbrace{\frac{1}{\bar\tau}\,z_k}_{\text{ceiling}} .
\]
\end{corollary}
\begin{remark}[The two ends answer to different things]\label{rem:asym}
Since $0<z_k<1$ we have $z_k^2<z_k$, so the bracket has a \emph{gap between exponents}:
unconditionally one can only conclude $\Rk\sim z_k^\alpha$ with $\alpha\in[1,2]$, hence
$k^\ast\in[\Theta(T^{1/3}),\Theta(T^{1/2})]$. Crucially the gap is asymmetric:
\begin{itemize}
\item the \textbf{floor} $z_k^2$ is produced by the \emph{dispersion} $V_\delta$
  (numerator of the Gram bound) and is present whenever the forecasters are
  heterogeneous --- coherence plays no role in the exponent;
\item the \textbf{ceiling} $z_k$ is the Jensen/average-subset bound, which is
  intrinsically first order because it cannot see that bagging cancels the leading term.
\end{itemize}
The effective exponent is therefore read off the \emph{ceiling}: the floor is always
$z_k^2$, and where $\alpha$ actually lands is decided by how far the ceiling can be pulled
down. This is the entire role of Tier~II: \emph{Assumption~A moves the ceiling,
Assumption~B sets the floor.}
\end{remark}

\begin{proposition}[Strict monotonicity]\label{prop:mono}
If each $g_i(k)$ is monotone in $k$, then $\Rk$ is strictly decreasing in $k$ for any
non-degenerate precision profile:
\[
  \frac{d\Rk}{dk}=\frac{2}{m^2}\sum_i\tau_i\Big(g_i-\frac1{\bar\tau}\Big)\dot g_i<0 .
\]
\end{proposition}
\begin{proof}
From Proposition~\ref{prop:identity} and $\sum_i\tau_i\dot g_i=0$ (differentiating the
conservation law), $\frac{d\Rk}{dk}=\frac2{m^2}\sum_i\tau_i(g_i-1/\bar\tau)\dot g_i$.
Each $g_i$ runs from $g_i(1)=1/\tau_i$ to $g_i(m)=1/\bar\tau$; under monotonicity, a
high-precision coordinate ($\tau_i>\bar\tau$) has $g_i\le1/\bar\tau$ throughout with
$\dot g_i\ge0$, and a low-precision one the reverse, so every product is $\le0$, strictly
for at least one $i$ unless all $\tau_i$ are equal.
\end{proof}
\begin{remark}
The monotonicity of each $g_i$ is the one lemma in Tier~I that requires its own proof;
the unconditional bound $\Rk\le R_1$ (Jensen on the harmonic mean) survives without it.
Proposition~\ref{prop:mono} has a sharp consequence: \emph{in population there is no
bias--variance trade-off in $k$} --- larger $k$ is always weakly better. An interior
optimum exists \emph{only} once estimation error enters through the finite-$T$ inflation
factor below.
\end{remark}

\begin{proposition}[Exact optimal-$k$ condition]\label{prop:foc}
Let $A=\sigma_\varepsilon^2+\Qopt$ and let the idealized infinite-bagged loss be
$L_k=(A+\Rk)\frac{T-1}{T-k}$. The discrete optimum $k^\ast$ satisfies
\[
  (T-k^\ast)\,(R_{k^\ast}-R_{k^\ast+1})=A+R_{k^\ast}.
\]
\end{proposition}

\begin{proposition}[RSM dominates CSR]\label{prop:csr}
For every $k$, $Q_k^{\mathrm{RSM}}\le Q_k^{\mathrm{CSR}}$: within each subset the
constrained-OLS weight minimizes $w'\Sig w$ over the equal-weight choice $1/k$; take
expectations over $S$.
\end{proposition}

%==========================================================================
\section{Tier II: the assumptions needed for a rate}
The unconditional bracket of Corollary~\ref{cor:open} is already two-sided, but its two
exponents differ: floor $z_k^2$, ceiling $z_k$. By Remark~\ref{rem:asym} the floor is
fixed by dispersion alone, so the only thing standing between us and a genuine rate
$\Rk\asymp z_k^2$ is closing the \emph{ceiling} from $z_k$ down to $z_k^2$. This requires
exactly two assumptions, and they are not interchangeable: Assumption~\ref{ass:A}
sharpens the ceiling, Assumption~\ref{ass:B} guarantees the floor is non-degenerate.

\begin{assumption}[Bounded coherence]\label{ass:A}
There exist constants $0<c\le C<\infty$, independent of $m$, such that
$c\le\tau_i/\bar\tau\le C$ for all $i$.
\end{assumption}
\noindent\emph{Role.} The upper part ($C$) is ``no dominant forecaster'':
$\kappa=\tau_{\max}/\bar\tau\le C$. The lower part ($c$) is ``no negligible
forecaster'': it forces $A_S\ge ck\bar\tau$, keeping $1/A_S$ and all its reciprocal
moments bounded. This is what makes the expansion
$\E_S[1/A_S]=\frac1{k\bar\tau}+\frac{\Var_S(A_S)}{(k\bar\tau)^3}+\cdots$ rigorous, with
an explicit, bounded remainder, rather than a heuristic.

\begin{assumption}[Non-degenerate dispersion]\label{ass:B}
There exists $v>0$, independent of $m$, such that $V_\delta\ge v$.
\end{assumption}
\noindent\emph{Role.} Genuine heterogeneity. Without it $\Rk\to0$ and there is nothing to
bound (simple averaging is near-optimal). The upper bound $V_\delta\le(C-1)^2$ and
$H\le C/c$ are already implied by Assumption~\ref{ass:A}.

\begin{theorem}[Two-sided $z_k^2$ bound]\label{thm:rate}
Under Assumptions~\ref{ass:A} and \ref{ass:B} there are constants $0<c_1\le c_2<\infty$,
depending only on $c,C,v$, such that for all $2\le k<m$,
\[
  \frac{c_1}{m\bar\tau}\,z_k^2\;\le\;\Rk\;\le\;\frac{c_2}{m\bar\tau}\,z_k^2 .
\]
Consequently the effective exponent is $\alpha=2$ exactly, and the optimal subset size
satisfies $k^\ast\asymp T^{1/3}$.
\end{theorem}
\begin{proof}[Proof sketch]
\emph{Upper.} Under Assumption~\ref{ass:A}, the reciprocal expansion of
$\E_R[1/L_{i,R}]$ has a remainder controlled by $A_S\ge ck\bar\tau$; combined with the
SRSWOR variance bound
$\Var_R(L_{i,R})\le\frac{(k-1)(m-k)}{m(m-1)(m-2)}V_\delta$ this gives
$|\bar v_{k,i}/\ell_i-1|\le K_{c}z_k(|h_i|+V_\delta)$ with $h_i=\tau_i/\bar\tau-1$.
Plugging into Proposition~\ref{prop:identity} and using $\ell_i\le C/m$,
$\sum_i\ell_ih_i^2\le CV_\delta$, yields
$\Rk\le\frac{K_{c,C}}{m\bar\tau}(V_\delta+V_\delta^2)z_k^2$. \emph{Lower.} In
Proposition~\ref{prop:bracket}, $k\E_S[1/A_S]-1/\bar\tau\asymp V_\delta z_k$ and the
denominator $\sum_i\sigma_i^2-m/\bar\tau=m\bar\tau^{-1}(H-1)$ is bounded by
Assumption~\ref{ass:A}; with $V_\delta\ge v$ this gives
$\Rk\ge\frac{v^2}{m(H-1)\bar\tau}z_k^2$. The rate for $k^\ast$ follows by inserting
$\Rk\asymp C_R z_k^2$ (with $z_k\approx1/k$ for $k\ll m$) into
Proposition~\ref{prop:foc}: $Ak^3+3C_Rk-2C_RT=0$, so $k^\ast\approx(2C_RT/A)^{1/3}$.
\end{proof}

\begin{remark}[How each assumption acts, and two caveats]\label{rem:caveats}
Relative to Corollary~\ref{cor:open}, the theorem changes \emph{only the ceiling}:
Assumption~\ref{ass:A} replaces the Jensen $z_k$ bound by a $z_k^2$ bound, collapsing the
bracket onto the floor that Assumption~\ref{ass:B} already pins at $z_k^2$. Two caveats
keep this honest.
\begin{enumerate}
\item \emph{The floor stays valid but loosens under high coherence.} If a dominant
forecaster does appear, the true $\Rk$ rises toward the $z_k$ ceiling ($\alpha\to1$), so
the $z_k^2$ floor remains a correct lower bound but is no longer tight. A floor never
certifies $\alpha=2$; only the ceiling coming down does.
\item \emph{The floor's constant is uniform in $m$ only under Assumption~\ref{ass:A}.}
The lower constant is $\propto V_\delta^2/\big(m(H-1)\bar\tau\big)$, and $H$ is bounded
uniformly in $m$ only under \ref{ass:A} (a near-negligible forecaster makes $H$ blow up
and the floor sag). For a fixed profile this is irrelevant --- $H$ is a finite number and
the floor is a fixed constant times $z_k^2$ --- so the $z_k^2$ \emph{exponent} on the
floor is owed to \ref{ass:B} alone; only its uniform constant borrows from \ref{ass:A}.
\end{enumerate}
\end{remark}

%==========================================================================
\section{The explicit leading constant}\label{sec:constant}
Theorem~\ref{thm:rate} shows $\Rk\asymp z_k^2$ with constants $c_1,c_2$ left abstract.
Under the same assumptions one can compute the leading constant in closed form from the
first three moments of the precision profile. The result instantiates $c_1=c_2=D_\tau$
(up to the $(m/(m-1))^2$ factor) \emph{inside} the bracket; it is a leading-order
approximation, not a new bound.

Recall $h_i=\tau_i/\bar\tau-1$ (so $\sum_ih_i=0$ and $V_\delta=\frac1m\sum_ih_i^2$), and
define the third moment and the \emph{weighted heterogeneity constant}
\[
  M_3=\frac1m\sum_{i=1}^m h_i^3,\qquad
  C_\tau=\frac1m\sum_{i=1}^m(1+h_i)(h_i-V_\delta)^2=V_\delta+M_3-V_\delta^2 .
\]
Since $1+h_i=\tau_i/\bar\tau=m\ell_i>0$, we have $C_\tau=\sum_i\ell_i(h_i-V_\delta)^2\ge0$:
it is exactly the variance of $h$ under the leverage (precision-weighted) distribution
$\ell_i=\tau_i/A_m$, linking it to the coherence picture. If heterogeneity is mild,
$M_3,V_\delta^2$ are higher order and $C_\tau\approx V_\delta$.

\begin{lemma}[Conservation-consistent first-order expansion]\label{lem:expand}
Under Assumption~\ref{ass:A}, to first order in $z_k$,
\[
  g_i(k)=\frac1{\bar\tau}\Big[1+\frac{m}{m-1}\,z_k\,(V_\delta-h_i)\Big]+o(z_k).
\]
\end{lemma}
\begin{proof}[Derivation]
Conditioning on $i\in S$, $\E_S[A_S\mid i\in S]=k\bar\tau\big[1+\frac{m}{m-1}z_k\,h_i\big]$,
so the delta method gives $g_i\approx\frac1{\bar\tau}[1-\frac{m}{m-1}z_k\,h_i]$. This
violates the exact identity $\sum_i\tau_ig_i=m$ (Lemma~\ref{lem:moments}) by
$-\frac{m}{m-1}z_k\,mV_\delta$. The discrepancy is removed by recentering $h_i$ at its
\emph{precision-weighted} mean, which is exactly $V_\delta$:
$\frac{\sum_i\tau_ih_i}{\sum_i\tau_i}=\frac1m\sum_i(h_i+h_i^2)=V_\delta$. Replacing $h_i$
by $h_i-V_\delta$ restores $\sum_i\tau_ig_i=m$ identically and yields the stated form; the
second-order reciprocal term contributes at $o(z_k)$ under Assumption~\ref{ass:A}.
\end{proof}

\begin{proposition}[Closed-form leading constant]\label{prop:Dtau}
Under Assumptions~\ref{ass:A}--\ref{ass:B}, with $B=1/\bar\tau$,
\[
  \Rk\approx D_\tau\,z_k^2,\qquad
  D_\tau=\frac{B}{m}\Big(\frac{m}{m-1}\Big)^2 C_\tau,
\]
equivalently $Q_k^{\mathrm{bag},\infty}\approx\Qopt+D_\tau z_k^2$ with $\Qopt=q+B/m$. In
the mild-heterogeneity limit $C_\tau\to V_\delta$, so $D_\tau\to V_\delta/(m\bar\tau)$,
recovering the H\'ajek/$\alpha=2$ leading term.
\end{proposition}
\begin{proof}
Substitute Lemma~\ref{lem:expand} into the exact identity
$\Rk=\frac1{m^2}\sum_i\tau_i(g_i-\frac1{\bar\tau})^2$. With
$g_i-\frac1{\bar\tau}=-\frac1{\bar\tau}\frac{m}{m-1}z_k(h_i-V_\delta)$,
\[
  \Rk\approx\frac1{m^2}\Big(\frac{m}{m-1}\Big)^2\frac{z_k^2}{\bar\tau^2}
  \sum_i\tau_i(h_i-V_\delta)^2
  =\frac1{m\bar\tau}\Big(\frac{m}{m-1}\Big)^2 C_\tau\,z_k^2,
\]
using $\sum_i\tau_i(h_i-V_\delta)^2=\bar\tau\sum_i(1+h_i)(h_i-V_\delta)^2=\bar\tau mC_\tau$.
\end{proof}

\begin{corollary}[Explicit cube-root rule]\label{cor:kstar-explicit}
With $A=\sigma_\varepsilon^2+\Qopt$ and $L_k=(A+D_\tau z_k^2)\frac{T-1}{T-k}$, the exact
first-order condition is $A+D_\tau z_k^2=\dfrac{2D_\tau z_k(T-k)}{k^2}$, which for large
$T$ (and $z_k\approx1/k$) gives
\[
  k^\ast\approx\Big(\frac{2D_\tau T}{A}\Big)^{1/3} .
\]
This is a directly computable plug-in: estimate $\bar\tau$, $V_\delta$, $M_3$, form
$D_\tau$, and minimize $L_k$ over integer $k$.
\end{corollary}

\begin{remark}[Status: approximation, not bound]\label{rem:status}
Proposition~\ref{prop:Dtau} is the \emph{leading-order} term, obtained from a first-order
expansion of $1/A_S$; it is not a two-sided bound and does not supersede
Proposition~\ref{prop:bracket}. Its role is to supply the explicit value of the constant
inside the closed bracket of Theorem~\ref{thm:rate}. Three qualifications:
(i) the expansion of $1/A_S$ requires $A_S$ bounded away from zero, i.e.\
Assumption~\ref{ass:A}; (ii) it is a single power of $z_k$, accurate for mild-to-moderate
heterogeneity --- for strongly heterogeneous (factor-type) profiles the curvature at
small $k$ is governed by $M_3$ and higher moments, so a single leading term understates
the small-$k$ decay; (iii) when the precisions $\tau_i$ are estimated, $D_\tau$ inherits
the sampling noise in $V_\delta$ and especially $M_3$. Within these limits it is the
correct and useful closed form for the $z_k^2$ constant, consistent in the mild limit
with the H\'ajek term and with the $T^{1/3}$ scaling observed in the Monte Carlo.
\end{remark}

%==========================================================================
\section{Convergence along the path}\label{sec:path}
Proposition~\ref{prop:endpoints} says $\bar v_k$ is a \emph{path} in weight space from
$w^{ave}=\frac1m\iotab$ (at $k=1$) to $w^\ast$ (at $k=m$), so
$\Rk=\|\bar v_k-w^\ast\|_\tau^2$ is the squared \emph{distance still to travel}, where
$\langle a,b\rangle_\tau=\sum_i a_ib_i/\tau_i$. Bounding $\Rk$ is then a question about how
the weights converge, and heterogeneity governs the shape of that convergence. Write the
per-coordinate fractional progress
\[
  \pi_i(k)=\frac{\bar v_{k,i}-1/m}{\ell_i-1/m}\in[0,1],\qquad \pi_i(1)=0,\ \ \pi_i(m)=1 .
\]

\begin{proposition}[Homogeneous convergence $\Rightarrow$ straight line]\label{prop:homog}
If the coordinates progress uniformly, $\pi_i(k)\equiv\pi(k)$, then the path is the line
$w^{ave}\!\to\! w^\ast$, the deviation $\delta_k=\bar v_k-w^\ast$ stays parallel (in the
$\tau$-metric) to $\eta=w^{ave}-w^\ast$, and
\[
  \Rk=(1-\pi(k))^2\,R_1,\qquad R_1=\frac{H-1}{m\bar\tau}.
\]
In this case the Gram/projection lower bound of Proposition~\ref{prop:bracket} holds with
\emph{equality}.
\end{proposition}
\noindent This recovers the tightness condition of Proposition~\ref{prop:bracket}
(equality for $\le2$ distinct precision types) in path language: uniform convergence
$\Leftrightarrow$ $\delta_k\parallel_\tau\eta$ $\Leftrightarrow$ the projection bound is exact.

\begin{remark}[Heterogeneity bends the path]\label{rem:bend}
With genuine heterogeneity the coordinates do not move in lockstep. From
Lemma~\ref{lem:expand}, to leading order
\[
  \bar v_{k,i}-w^\ast_i\approx z_k\,u_i,\qquad u_i\propto \tau_i(h_i-V_\delta)
  =\bar\tau(1+h_i)(h_i-V_\delta),
\]
so convergence is self-similar in $z_k$ but along the direction $u$, whereas the
projection of Proposition~\ref{prop:bracket} uses the single direction
$\eta_i=\tfrac1m-\ell_i=-h_i/m$. The two coincide only in the homogeneous/$\le2$-type
limit; in general the \emph{angle between $u$ and $\eta$ is the heterogeneity}. This is
exactly why the single-direction bound is loose:
\begin{itemize}
\item projecting onto $\eta$ keeps only $u$'s $\eta$-component, giving the floor constant
  $V_\delta^2/(H-1)$ (Cor.~\ref{cor:open});
\item using the true direction $u$ keeps all of it, giving back $\Rk\approx D_\tau z_k^2$
  with $C_\tau=\frac1m\sum_i(1+h_i)(h_i-V_\delta)^2$ (Prop.~\ref{prop:Dtau}).
\end{itemize}
The two agree in the mild limit, where $H-1\to V_\delta$ and $C_\tau\to V_\delta$.
\end{remark}

\begin{remark}[The exponent is the slowest coordinate's convergence rate]\label{rem:slowest}
Since $\Rk=\sum_i\delta_{k,i}^2/\tau_i$, the exponent $\alpha$ is set by which weight
converges slowest.
\begin{itemize}
\item Bounded coherence (heterogeneous, no dominant forecaster): every coordinate
  converges at the common rate $\delta_{k,i}\sim z_k$, so $\Rk\sim z_k^2$ and $\alpha=2$.
\item Dominant forecaster (high coherence): a single coordinate $i^\ast$ with huge
  $\tau_{i^\ast}$ is heavily under-weighted at equal weights ($1/m\ll\ell_{i^\ast}$) and
  converges slowly, retaining a persistent $\eta$-aligned component; its term dominates
  and $\Rk\sim z_k$, so $\alpha\to1$.
\end{itemize}
Thus the bagging exponent is governed by the convergence rate of the slowest-converging
weight --- the convergence-theoretic face of the floor/ceiling bracket.
\end{remark}

\noindent\textbf{Usable form.} The path view gives the anchored distance-remaining model
\[
  \Rk=R_1\,\phi(z_k)^2,\qquad \phi(z_1)=1,\ \ \phi(0)=0,\ \
  \phi(z_k)\approx\sqrt{D_\tau/R_1}\;z_k\ \ (\text{leading order}),
\]
with $\phi$ the normalized distance still to travel and heterogeneity setting its shape
through $u$: homogeneous convergence gives $\phi(z_k)=z_k/z_1$ ($\alpha=2$), a dominant
forecaster flattens $\phi$ toward $\sqrt{z_k}$ ($\alpha\to1$), and minimization returns
$k^\ast\propto T^{1/(\alpha+1)}$ as in Corollary~\ref{cor:kstar-explicit}. To tighten the
\emph{lower bound} rigorously, project onto the two-dimensional span $\{\eta,u\}$
(equivalently append the moment $\sum_i\tau_i^2g_i$): exact for $\le3$ precision types, the
geometric form of the moment hierarchy of Proposition~\ref{prop:bracket}.

%==========================================================================
\section{U-statistic structure and finite-ensemble convergence}\label{sec:ustat}
Two idealizations sit outside the bracket: $k$ is treated as a growing tuning size, and
bagging is taken to the infinite-ensemble limit $G\to\infty$. The generalized-U-statistic
theory of Peng, Coleman \& Mentch (2022) closes the first, and the algorithmic-convergence
analysis of Lopes, Wu \& Lee (2019) the second.

\subsection*{The bagged weight is a generalized U-statistic}
The infinite-bagged weight $\bar v_{k,i}=\E_S[h_i(S)]$ is a complete generalized
U-statistic of degree $k$ on the $m$ forecasters, with kernel $h_i(S)=\tau_i/A_S$; the
realized RSM weight $\hat{\bar v}_{k,i}=\frac1G\sum_{g=1}^G v_{S_g,i}$ over $G$ random
subsets is its incomplete version. Write $\zeta_k=\Var_S(h_i)$ (full kernel variance) and
$\zeta_1=\Var\big(\E[h_i\mid\text{one forecaster fixed}]\big)$ (the H\'ajek/first-order
variance). The variance ratio $\zeta_k/(k\zeta_1)$ measures asymptotic linearity: it tends
to $1$ exactly when the H\'ajek projection dominates.

\begin{remark}[H\'ajek projection is the exact leading term]\label{rem:hajek}
Because $h_i=\tau_i/A_S$ is a smooth, OLS-type functional of the subset precisions, its
variance ratio is bounded (Peng--Coleman--Mentch, Prop.~1--4 and App.~B.3 for OLS), so
$\bar v_{k,i}$ is asymptotically linear and the H\'ajek projection \emph{is} the leading
term. This is the rigorous justification for treating the $z_k^2$ expansion of
Proposition~\ref{prop:Dtau} as the exact leading order (not a heuristic): the
small-heterogeneity result and the H\'ajek first-order projection coincide, fixing
$\alpha=2$ under bounded coherence. The Berry--Esseen bound of Peng--Coleman--Mentch
(Thm.~3) quantifies the rate of normality as $\propto n^{-1/2}+\{\tfrac kn(\tfrac{\zeta_k}
{k\zeta_1}-1)\}^{1/2}$, making the roles of $k$ and the sample size explicit.
\end{remark}

\subsection*{Finite-ensemble ($G<\infty$) inflation}
Under common loading the gap between the realized finite-$G$ risk and the infinite-bagged
risk is exact and of order $1/G$.

\begin{proposition}[Finite-$G$ risk inflation]\label{prop:finiteG}
Let $\hat{\bar v}_k=\frac1G\sum_{g=1}^G v_{S_g}$ with $S_1,\dots,S_G$ i.i.d.\ uniform
size-$k$ subsets, and $\hat Q_G=\hat{\bar v}_k'\Sig\hat{\bar v}_k$. Since each $v_{S_g}$
sums to one, so does $\hat{\bar v}_k$, and
\[
  \E_G[\hat Q_G]=Q_k^{\mathrm{bag},\infty}+\frac{\Gamma_k}{G},\qquad
  \Gamma_k=\sum_{i=1}^m\frac{\Var_S(v_{S,i})}{\tau_i}\ \ge0 .
\]
\end{proposition}
\begin{proof}
The draws are i.i.d., so $\E_G[\hat{\bar v}_{k,i}]=\bar v_{k,i}$ and
$\Var_G(\hat{\bar v}_{k,i})=\Var_S(v_{S,i})/G$. As in
Proposition~\ref{prop:identity} the common factor and cross terms drop for a sum-to-one
vector, giving $\hat Q_G=q+\sum_i\hat{\bar v}_{k,i}^2/\tau_i$; take expectations and use
$\E_G[\hat{\bar v}_{k,i}^2]=\bar v_{k,i}^2+\Var_S(v_{S,i})/G$.
\end{proof}

\begin{remark}[Reading $\Gamma_k$ and choosing $G$]\label{rem:G}
$\Gamma_k$ is the ensemble (Monte-Carlo) variance --- the common-loading analogue of the
Peng--Coleman--Mentch incomplete-U-statistic term $\zeta_k/N$ with $N=G$, and it matches
the $O(1/t)$ algorithmic-convergence rate of Scornet (2016) and Lopes--Wu--Lee. It
\emph{decreases in $k$}: at $k=1$, $v_{S,i}$ is Bernoulli$(1/m)$ with variance
$\tfrac1m(1-\tfrac1m)$, while at $k=m$ it is the deterministic $\ell_i$ with zero variance.
Hence \textbf{small $k$ needs more bags}. Practically, require the finite-$G$ inflation
$\Gamma_k/G$ to be small relative to the bagging gain $R_1-R_k$; or, without computing
$\Gamma_k$, use the Lopes--Wu--Lee \emph{ensemble bootstrap} --- resample the $G$
subset-weights with replacement to estimate the $(1-\alpha)$-quantile $q_{1-\alpha}(G)$ of
$\hat Q_G-Q_k^{\mathrm{bag},\infty}$ from a single run --- and raise $G$ until
$q_{1-\alpha}(G)$ is below tolerance. Lopes--Wu--Lee (\S2.2) also note that some
coordinates' weights converge more slowly, distorting them: the same slowest-coordinate
effect as Remark~\ref{rem:slowest}.
\end{remark}

\begin{remark}[Two error sources in the endogenous case]\label{rem:endoG}
When the weights are estimated from $T$ observations there are two independent error
sources, and the incomplete-U-statistic variance takes the Peng--Coleman--Mentch form
\[
  \approx \underbrace{\frac{C_1}{T}}_{\text{estimation (data)}}
  +\underbrace{\frac{\Gamma_k}{G}}_{\text{ensemble}} .
\]
The variance ratio $\zeta_k/(k\zeta_1)$ governs which dominates; the two are comparable
when $G\sim T\,\Gamma_k/C_1$, so $G\sim T$ suffices and $G\gg T$ is unnecessary. The data
term is the estimation-error/Kan--Zhou contribution that drives the gap between the
exogenous and endogenous optimal $k$.
\end{remark}

%==========================================================================
\section{Minimality: each failure mode collapses one side}
The two assumptions are tight; because each governs a different end of the bracket,
dropping either one degrades the bound in a specific, interpretable way. The following
table reads off the floor and ceiling under each combination.

\begin{center}
\begin{tabular}{lcccc}
\hline
Assumptions & Floor & Ceiling & $\alpha$ & $k^\ast$ \\
\hline
\ref{ass:A} and \ref{ass:B} & $z_k^2$ & $z_k^2$ & $2$ & $\Theta(T^{1/3})$ \\
\ref{ass:B} only (no \ref{ass:A}) & $z_k^2$ & $z_k$ & $[1,2]$ & $[\Theta(T^{1/3}),\Theta(T^{1/2})]$ \\
\ref{ass:A} only (no \ref{ass:B}) & $0$ & $z_k^2$ & $[0,2]$ & up to $\Theta(T^{1/3})$ \\
neither & $0$ & $z_k$ & $[0,2]$ & --- \\
\hline
\end{tabular}
\end{center}

\noindent The middle row is the case discussed throughout: drop bounded coherence and the
\emph{floor is unchanged} at $z_k^2$ (it is owed to dispersion, \ref{ass:B}), while the
\emph{ceiling springs back} from $z_k^2$ to the Jensen value $z_k$ --- precisely the open
bracket of Corollary~\ref{cor:open}. The economic content of each failure mode:

\begin{center}
\begin{tabular}{lll l}
\hline
Assumption dropped & Regime & What survives & Rate \\
\hline
\ref{ass:A}-lower ($\tau_{\min}/\bar\tau\to0$) & dominant/negligible forecaster
  & only Jensen, $\Rk\lesssim z_k$ & $k^\ast\asymp T^{1/2}$ \\
\ref{ass:B}-lower ($V_\delta\to0$) & homogeneous forecasters
  & $\Rk\to0$ & no tuning; averaging optimal \\
\hline
\end{tabular}
\end{center}

\noindent Thus the cube root is precisely the \emph{bounded-coherence, heterogeneous}
regime; the square root and the degenerate case are its two boundaries. The effective
exponent $\alpha\in[1,2]$ is therefore not a fitted quantity but a map of which
assumption is binding --- and, by Remark~\ref{rem:asym}, it is the ceiling that moves.

%==========================================================================
\section{Discussion: plausibility, and a fixed-$m$ alternative}
\paragraph{Empirical content.} The two assumptions are checkable, and they align with the
applications. Assumption~\ref{ass:B} (heterogeneity) is the ``there are gains'' condition:
it holds in the heterogeneous vintage panels, where RSM delivers real gains, and fails for
near-homogeneous professional-forecaster consensus data, where gains are negligible ---
exactly the dividing line the empirical work traces. Assumption~\ref{ass:A} (no dominant
forecaster) is the same economic point made in the forecast-combination literature: a
single super-precise or strongly negatively correlated forecaster is the case where
averaging fails badly and is empirically unlikely; assuming \ref{ass:A} assumes that
pathology away.

\paragraph{Uniformity is only for the large-$m$ story.} Assumptions~\ref{ass:A} and
\ref{ass:B} are stated with constants independent of $m$ because that is what a
\emph{rate in $m$} (a stable exponent as the panel grows) requires. If only a fixed panel
is of interest, no uniformity is needed: the exact bracket
(Proposition~\ref{prop:bracket}) already supplies finite, computable two-sided bounds, and
the only condition required for a non-trivial statement is the non-degeneracy
$V_\delta>0$. In short:
\[
  \text{fixed }m:\ \ V_\delta>0 \text{ plus the exact bracket};\qquad
  \text{large }m / \text{rate}:\ \ \text{uniform Assumptions~\ref{ass:A} and \ref{ass:B}.}
\]

\paragraph{Scope.} All of the above governs the population object $\Rk$ (infinite
bagging). The map from $\Rk$ to a realized, finite-$G$, finite-$T$ mean-squared forecast
error still involves the bagging idealization ($G\to\infty$) and the estimation-error
inflation factor $\frac{T-1}{T-k}$, which are modeling choices outside the bracket and
should be stated as such.

%==========================================================================
\section{Summary}
\begin{enumerate}
\item \textbf{Exact, no assumptions:} oracle weights are leverage scores
(Prop.~\ref{prop:oracle}); $\Rk$ is an exact quadratic form in the weight deviation
(Prop.~\ref{prop:identity}); RSM bridges averaging ($k=1$) and the oracle ($k=m$)
(Prop.~\ref{prop:endpoints}); a computable two-sided bracket
(Prop.~\ref{prop:bracket}); strict monotonicity, so no population bias--variance
trade-off (Prop.~\ref{prop:mono}); an exact optimal-$k$ condition
(Prop.~\ref{prop:foc}); and RSM $\succeq$ CSR (Prop.~\ref{prop:csr}).
\item \textbf{Unconditional open bracket:} $z_k^2\lesssim\Rk\lesssim z_k$
(Cor.~\ref{cor:open}), with the floor $z_k^2$ set by dispersion and the ceiling $z_k$ the
Jensen bound; hence $\alpha\in[1,2]$, $k^\ast\in[\Theta(T^{1/3}),\Theta(T^{1/2})]$.
\item \textbf{Rate, two assumptions:} \emph{Assumption~\ref{ass:A} moves the ceiling}
($z_k\to z_k^2$); \emph{Assumption~\ref{ass:B} sets the floor} ($z_k^2$). Together the
bracket closes: $\Rk\asymp z_k^2$ and $k^\ast\asymp T^{1/3}$ (Thm.~\ref{thm:rate}).
\item \textbf{Explicit constant:} under both assumptions the leading constant is closed
form, $\Rk\approx D_\tau z_k^2$ with $D_\tau=\frac{B}{m}(\frac{m}{m-1})^2C_\tau$ and
$C_\tau=V_\delta+M_3-V_\delta^2$ the leverage-weighted dispersion of $h$
(Prop.~\ref{prop:Dtau}); it gives the plug-in $k^\ast\approx(2D_\tau T/A)^{1/3}$
(Cor.~\ref{cor:kstar-explicit}) and reduces to the H\'ajek term $V_\delta/(m\bar\tau)$ in
the mild limit. It is a leading-order approximation inside the bracket, not a new bound
(Rem.~\ref{rem:status}).
\item \textbf{Convergence view:} $\Rk=\|\bar v_k-w^\ast\|_\tau^2$ is the squared
distance-to-go on the averaging$\to$oracle path (Prop.~\ref{prop:endpoints}). Uniform
convergence gives $\Rk=(1-\pi(k))^2R_1$ and makes the projection bound exact
(Prop.~\ref{prop:homog}); heterogeneity bends the convergence direction from $\eta$ toward
$u\propto\tau(h-V_\delta)$ (Rem.~\ref{rem:bend}), and $\alpha$ is the slowest coordinate's
convergence rate (Rem.~\ref{rem:slowest}).
\item \textbf{U-statistic / finite ensemble:} $\bar v_{k,i}$ is a generalized U-statistic
with kernel $\tau_i/A_S$; its bounded variance ratio makes the H\'ajek projection the exact
leading term, rigorizing $\alpha=2$ (Rem.~\ref{rem:hajek}). The realized finite-$G$ risk
satisfies $\E_G[\hat Q_G]=Q_k^{\mathrm{bag},\infty}+\Gamma_k/G$
(Prop.~\ref{prop:finiteG}), with $\Gamma_k$ decreasing in $k$ (small $k$ needs more bags);
in the endogenous case the variance splits as $C_1/T+\Gamma_k/G$, so $G\sim T$ suffices
(Rems.~\ref{rem:G}--\ref{rem:endoG}).
\item \textbf{Minimality:} dropping \ref{ass:A} reopens the ceiling to $z_k$
(square-root regime, $\alpha\to1$); dropping \ref{ass:B} drops the floor to $0$
(degenerate, averaging-optimal). The exponent $\alpha\in[1,2]$ indexes which assumption
binds, and it is the ceiling that moves.
\end{enumerate}

\end{document}
